\documentclass[11pt,a4paper]{article}
\newcommand{\DocVersion}{1.1.1}
\newcommand{\BuildDate}{20261026}
% --------------------------------------------------
% Packages
% --------------------------------------------------
\usepackage[utf8]{inputenc}
\usepackage[T1]{fontenc}
\usepackage[english]{babel}

\usepackage{geometry}
\usepackage{xcolor}
\usepackage{graphicx}
\usepackage{enumitem}
\usepackage{hyperref}
\usepackage{fancyhdr}
\usepackage{booktabs}
\usepackage{tabularx}
\usepackage{array}
\usepackage{longtable}

% --------------------------------------------------
% Page settings
% --------------------------------------------------
\geometry{
    top=2.5cm,
    bottom=2.5cm,
    left=2.5cm,
    right=2.5cm
}

\definecolor{companyblue}{HTML}{1F4E79}
\definecolor{lightblue}{HTML}{EAF2F8}
\definecolor{warningred}{HTML}{9C0006}
\definecolor{warningbg}{HTML}{FFC7CE}

\hypersetup{
    colorlinks=true,
    linkcolor=companyblue,
    urlcolor=companyblue
}

\setlength{\parindent}{0pt}
\setlength{\parskip}{0.6em}

% --------------------------------------------------
% Header and footer
% --------------------------------------------------
\pagestyle{fancy}
\fancyhf{}

\lhead{Tex C. Lab}
\rhead{User Manual}
\cfoot{\thepage}

% --------------------------------------------------
% Custom commands
% --------------------------------------------------
\newcommand{\software}{\textbf{Tex C. Lab}}
\newcommand{\button}[1]{\textbf{#1}}
\newcommand{\filename}[1]{\texttt{#1}}

% --------------------------------------------------
% Document
% --------------------------------------------------
\begin{document}

% ==================================================
% Cover page
% ==================================================
\begin{titlepage}
    \centering

    \vspace*{1.5cm}

    % Company logo
    % The file must be located in Docs/imgs/logo.png
    \includegraphics[width=0.35\textwidth]{imgs/cra-logo.png}

    \vspace{1.5cm}

    {\Huge\bfseries Tex C. Lab\par}

    \vspace{0.5cm}

    {\Large Utility for managing encrypted files for Power systems\par}

    \vspace{2cm}

    {\LARGE\bfseries User Manual\par}

    \vspace{1cm}

    {\large Version \DocVersion\par}

    \vspace{0.3cm}

    {\large Build date: \BuildDate\par}

    \vfill

    {\large Tex \par}

\end{titlepage}

% ==================================================
% Table of contents
% ==================================================
\tableofcontents
\newpage

% ==================================================
\section{Introduction}
% ==================================================

\software{} is a software tool developed to generate encrypted
files used by Power systems and their Boot Manager.

The program can generate the files required to:

\begin{itemize}
    \item Automatically log in to the Boot Manager;
    \item Change the Administrator user's password;
    \item Reset the Administrator user's password;
    \item Update the card parameters;
    \item Validate new firmware for the card;
    \item Generate a temporary authorization file granting
    Engineer permissions to a user.
\end{itemize}

Depending on the operation, the generated files must either be
transferred to the card using a USB device or copied directly
to a specific directory on the controller.

This manual describes the features available in the application
and the procedures to follow when using the generated files
with the Boot Manager or the controller.

% ==================================================
\section{Requirements and general operation}
% ==================================================

The Boot Manager manages two access levels:

\begin{itemize}
    \item \textbf{Guest};
    \item \textbf{Administrator}.
\end{itemize}

Automatic login and operations related to the Administrator user's
password are managed through files generated by \software.

The application also provides a generator for temporary Engineer
authorization. This operation follows a different procedure:
the generated file must be copied to the controller's
\filename{F:/public} directory, and the controller must then
be rebooted.

Unlike the files used for automatic login and password management,
the temporary authorization file is not processed through the
USB/pre-bootloader workflow described below.

\subsection{Standard file processing procedure}

For login, password management, card parameter, and firmware
operations, where applicable, the generated files must be
transferred to the card using a USB device.

To process files through the pre-bootloader, it is necessary to:

\begin{enumerate}
    \item Copy the generated file to a USB device;
    \item Connect the USB device to the card;
    \item Power on or restart the card;
    \item Press and hold the \texttt{X} button during startup;
    \item Enter the \textbf{pre-bootloader};
    \item Select the required operation.
\end{enumerate}

\textbf{Important:} the Temporary Engineer Authorization file
does not use this procedure. Refer to the
\button{Temporary Engineer Authorization} section for its
specific workflow.

\subsection{Pre-bootloader}

The pre-bootloader is the startup environment accessed by holding
the \texttt{X} button while powering on the card.

From this environment, it is possible to:

\begin{itemize}
    \item Access the Boot Manager main menu;
    \item Load the card parameters;
    \item Load new firmware.
\end{itemize}

To process the generated files through the pre-bootloader, it is
necessary to access this environment. For login, change, and reset
operations, it is sufficient to return to the main menu. For
parameters and firmware, the corresponding options must instead
be selected.

Temporary Engineer Authorization is an exception to this workflow
and is handled directly by the controller.

% ==================================================
\section{Application interface}
% ==================================================

The application consists of the following pages:

\begin{itemize}
    \item \textbf{Manual};
    \item \textbf{About};
    \item \textbf{Settings};
    \item \textbf{Generators}.
\end{itemize}

% ==================================================
\subsection{Manual page}
% ==================================================

The Manual page redirects to a PDF viewer and automatically opens
the user manual.

% ==================================================
\subsection{About page}
% ==================================================

The \textbf{About} page contains information about the application,
including:

\begin{itemize}
    \item Software name;
    \item Installed version;
    \item Information about the software manufacturer.
\end{itemize}

% ==================================================
\subsection{Settings page}
% ==================================================

The \textbf{Settings} page allows you to select the destination
directory for files generated by the application.

\subsubsection*{Selecting the output directory}

To set the destination directory:

\begin{enumerate}
    \item Open the \button{Settings} page;
    \item Select the desired directory;
    \item Confirm the selection.
\end{enumerate}

All subsequently generated files will be saved in the configured
directory.

% ==================================================
\subsection{Generators page}
% ==================================================

The \textbf{Generators} page contains the generators used to prepare
files for use with the Boot Manager and the controller.

The available generators are:

\begin{itemize}
    \item \button{Generate Login};
    \item \button{Change Password};
    \item \button{Reset Password};
    \item \button{Card Parameters};
    \item \button{Firmware Checksum};
    \item \button{Temporary Engineer Authorization}.
\end{itemize}

% ==================================================
\section{Managing generated files}
% ==================================================

The files generated by \software{} are divided into three categories
according to how they are processed.

\subsection{Automatically processed files}

The following files are automatically processed by the pre-bootloader
when they are present on the USB device:

\begin{itemize}
    \item Login file;
    \item Password change file;
    \item Password reset file.
\end{itemize}

For this reason, these files must be removed from the USB device
after use, when they are no longer needed.

\textbf{Warning:} keeping these files on a USB device connected to
the card could cause the associated operation to be executed
automatically during a subsequent pre-bootloader startup.

\subsection{Manually loaded files}

The following files are manually selected by the user directly
from the pre-bootloader:

\begin{itemize}
    \item Card parameter file;
    \item Firmware file and associated checksum file.
\end{itemize}

These files are not processed automatically and can therefore remain
on the USB device without causing the operation to be executed
unintentionally.

\subsection{Temporary authorization file}

The Temporary Engineer Authorization file follows a different
procedure from the files described above.

It must be copied directly to the controller's
\filename{F:/public} directory. The controller must then be rebooted
to apply the temporary Engineer authorization.

This file is not loaded from a USB device and does not require
access to the pre-bootloader.

\textbf{CRITICAL WARNING:} once the controller has restarted and
the authorization has been confirmed, the generated file must be
removed from \filename{F:/public} immediately.

Do not leave the authorization file in the directory after use.
Removing the file is a mandatory cleanup step and does not, by
itself, imply that previously granted permissions have been revoked.

% ==================================================
\section{Generators}
% ==================================================

The generators produce files that are processed by the Boot Manager
or the controller according to the operation being performed.

For operations handled by the pre-bootloader, the bootloader is
accessed by pressing and holding the controller button during startup.

Temporary Engineer Authorization is an exception: its file is
copied to \filename{F:/public}, and the controller is rebooted
without following the USB/pre-bootloader procedure.

% ==================================================
\subsection{Generate Login}
% ==================================================

The \button{Generate Login} generator creates a file that allows
automatic login to the Boot Manager using the Administrator
credentials.

\subsubsection*{Required data}

\begin{itemize}
    \item Current Administrator password.
\end{itemize}

\subsubsection*{Procedure}

\begin{enumerate}
    \item Open the \button{Generators} page.
    \item Select \button{Generate Login}.
    \item Enter the Administrator password.
    \item Start generating the file.
    \item Copy the generated file to a USB device.
    \item Connect the USB device to the card.
    \item Power on or restart the card while holding down
    the \texttt{X} button.
    \item Enter the pre-bootloader.
    \item Select access to the \textbf{Main Menu}.
    \item Allow the Boot Manager to process the login file.
\end{enumerate}

% ==================================================
\subsection{Change Password}
% ==================================================

The \button{Change Password} generator creates a file that allows
the Administrator password to be changed.

To perform this operation, both the current login file and the new
password are required.

\subsubsection*{Required data}

\begin{itemize}
    \item New Administrator password.
\end{itemize}

\subsubsection*{Procedure}

\begin{enumerate}
    \item Open the \button{Generators} page.
    \item Select \button{Change Password}.
    \item Select the current login file.
    \item Enter the new password.
    \item Start generating the file.
    \item Copy the generated file to a USB device.
    \item Also copy the current login file.
    \item Connect the USB device to the card.
    \item Power on or restart the card while holding down
    the \texttt{X} button.
    \item Enter the pre-bootloader.
    \item Select the \textbf{Main Menu}.
    \item Allow the Boot Manager to process the generated file.
\end{enumerate}

After use, remove both the login file and the password change file
from the USB device. If you want to continue using a login file,
you must regenerate it using the new password that was just set.

% ==================================================
\subsection{Reset Password}
% ==================================================

The \button{Reset Password} generator creates a temporary file
for resetting the Administrator password.

The reset file is sufficient on its own: it is not necessary to add
a login file or a password change file.

\subsubsection*{Required data}

\begin{itemize}
    \item MAC address of the card;
    \item Expiration date and time of the file.
\end{itemize}

The MAC address must be entered in one of the formats supported
by the application, for example:

\begin{itemize}
    \item \texttt{00:11:22:33:44:55};
    \item \texttt{00-11-22-33-44-55};
    \item \texttt{001122334455}.
\end{itemize}

The expiration date and time must be set in the future. The file
contains a Unix timestamp used by the Boot Manager to verify
its validity.

\subsubsection*{Procedure}

\begin{enumerate}
    \item Open the \button{Generators} page.
    \item Select \button{Reset Password}.
    \item Enter the correct MAC address of the card.
    \item Set a future expiration date and time.
    \item Start generating the file.
    \item Copy the generated file to a USB device.
    \item Connect the USB device to the card.
    \item Power on or restart the card while holding down
    the \texttt{X} button.
    \item Enter the pre-bootloader.
    \item Select the \textbf{Main Menu}.
    \item Allow the Boot Manager to process the reset file.
\end{enumerate}

The file cannot be used after the specified expiration time
and must be removed from the USB device after use.

% ==================================================
\subsection{Card Parameters}
% ==================================================

The \button{Card Parameters} generator prepares the card parameter
file for loading through the pre-bootloader.

The parameters must first be downloaded from the Boot Manager.
The resulting file, named \filename{POWER.CRD}, can be edited
as a text file.

After editing, the file must be processed by \software{}.
The program generates the encrypted file:

\begin{center}
    \filename{CARDPARS.ENC}
\end{center}

\subsubsection*{Procedure}

\begin{enumerate}
    \item Download the card parameters from the Boot Manager.
    \item Open the \filename{POWER.CRD} file with a text editor.
    \item Modify the required parameters.
    \item Save the modified file.
    \item Open the \button{Generators} page.
    \item Select \button{Card Parameters}.
    \item Select the \filename{POWER.CRD} file.
    \item Start generating the file.
    \item Copy \filename{CARDPARS.ENC} to a USB device.
    \item Connect the USB device to the card.
    \item Power on or restart the card while holding down
    the \texttt{X} button.
    \item From the pre-bootloader, select
    \button{Upload Card Parameters}.
    \item Select or confirm the \filename{CARDPARS.ENC} file.
    \item Wait for the operation to complete.
\end{enumerate}

The \filename{CARDPARS.ENC} file is manually selected by the
pre-bootloader and is not processed automatically. Therefore,
if desired, it can be left on the USB device.

% ==================================================
\subsection{Firmware Checksum}
% ==================================================

The \button{Firmware Checksum} generator validates new firmware
for loading onto the card by generating the checksum required
by the card.

The firmware must first be processed by \software{}
before it can be loaded by the Boot Manager.

The program generates the encrypted file:

\begin{center}
    \filename{CHECKSUM.ENC}
\end{center}

\subsubsection*{Required data}

\begin{itemize}
    \item Firmware binary file.
\end{itemize}

\subsubsection*{Procedure}

\begin{enumerate}
    \item Open the \button{Generators} page.
    \item Select \button{Firmware Checksum}.
    \item Select the firmware binary file.
    \item Start generating the file.
    \item Copy \filename{CHECKSUM.ENC} to a USB device
    together with the firmware itself.
    \item Connect the USB device to the card.
    \item Power on or restart the card while holding down
    the \texttt{X} button.
    \item From the pre-bootloader, select
    \button{Upload Firmware}.
    \item Wait for the update to complete.
\end{enumerate}

The \filename{CHECKSUM.ENC} file is manually selected by the
pre-bootloader and is not processed automatically.

% ==================================================
\subsection{Temporary Engineer Authorization}
% ==================================================

The \button{Temporary Engineer Authorization} generator creates
a file used to grant temporary Engineer permissions to a user.

Unlike the other generators, this file is not processed through
the USB device or the pre-bootloader. Instead, it must be copied
directly to the controller's \filename{F:/public} directory.

The controller must then be rebooted for the authorization
procedure to take effect.

\subsubsection*{Procedure}

\begin{enumerate}
    \item Open the \button{Generators} page.
    \item Select \button{Temporary Engineer Authorization}.
    \item Start generating the file.
    \item Locate the generated file in the configured output
    directory.
    \item Copy the generated file to the
    \filename{F:/public} directory on the controller.
    \item Verify that the file is present in the correct directory.
    \item Reboot the controller.
    \item After the controller has restarted, verify that the
    intended user has been granted temporary Engineer permissions.
    \item Immediately remove the generated authorization file
    from \filename{F:/public}.
    \item Verify that the file is no longer present in the directory.
\end{enumerate}

\textbf{CRITICAL WARNING:} removing the authorization file after
the reboot and successful authorization is mandatory. Do not leave
the file in \filename{F:/public} after use.

The file is intended for temporary authorization and must be
handled as a security-sensitive artifact. Generate and use it
only when authorized, and do not distribute or retain unnecessary
copies.

Removing the file is a required cleanup measure. The effect of
file removal on permissions that have already been granted depends
on the controller's authorization implementation; do not assume
that deleting the file automatically revokes those permissions.

% ==================================================
\section{Generator summary}
% ==================================================

\renewcommand{\arraystretch}{1.4}

\begin{longtable}{p{3.2cm} p{4.2cm} p{5.2cm}}
\toprule
\textbf{Generator} &
\textbf{Input} &
\textbf{Result and usage} \\
\midrule
\endhead

\button{Generate Login} &
Administrator password &
Automatic login file to be used from the Main Menu. \\

\button{Change Password} &
Current login file and new password &
File for changing the Administrator password from the Main Menu. \\

\button{Reset Password} &
MAC address and expiration &
Temporary file for resetting the password from the Main Menu. \\

\button{Card Parameters} &
Modified \filename{POWER.CRD} file &
Generates \filename{CARDPARS.ENC}, to be used with
\button{Upload Card Parameters}. \\

\button{Firmware Checksum} &
Firmware binary file &
Generates \filename{CHECKSUM.ENC}, to be used with
\button{Upload Firmware}. \\

\button{Temporary Engineer Authorization} &
No additional input specified &
Generates a temporary authorization file to be copied to
\filename{F:/public}. Reboot the controller to apply the
authorization, then remove the file immediately. \\

\bottomrule
\end{longtable}

% ==================================================
\section{Warnings}
% ==================================================

\begin{itemize}
    \item Always verify that the USB device contains the
    correct file before connecting it to the card.

    \item Remove the password change and reset files after use,
    as they may be processed automatically by the pre-bootloader.

    \item After changing the password, if you want to use automatic
    login, remember to regenerate the login file using the new
    password that was just set.

    \item Carefully verify the MAC address before generating
    a reset file.

    \item Make sure the reset file's expiration date and
    time are in the future.

    \item Do not manually modify encrypted files generated
    by the application.

    \item Do not interrupt power to the card during parameter
    or firmware updates.

    \item Use only firmware compatible with the card model.

    \item Generate and use Temporary Engineer Authorization
    files only when authorized.

    \item For Temporary Engineer Authorization, copy the generated
    file to \filename{F:/public} on the controller and reboot
    the controller to apply the authorization.

    \item Immediately after confirming that temporary Engineer
    authorization has been granted, remove the authorization
    file from \filename{F:/public}. Never leave the file in
    the directory after use.

    \item Do not assume that removing the authorization file
    automatically revokes permissions already granted by
    the controller.
\end{itemize}

% ==================================================
\section{Troubleshooting}
% ==================================================

\subsection{The file is not processed}

For files processed through the USB/pre-bootloader procedure,
check the following:

\begin{itemize}
    \item The file is located in the correct directory on the
    USB device;
    \item The file name has not been changed;
    \item The USB device was connected before startup;
    \item The \texttt{X} button was held down during power-on;
    \item The correct option was selected in the Boot Manager;
    \item The file has not expired;
    \item The file was generated for the correct card.
\end{itemize}

For Temporary Engineer Authorization, do not use the USB
or pre-bootloader procedure. Verify that the file was copied
to \filename{F:/public} on the controller and that the controller
was rebooted after the file was placed in the directory.

\subsection{Temporary Engineer authorization is not granted}

Check the following:

\begin{itemize}
    \item The correct file was generated using
    \button{Temporary Engineer Authorization};
    \item The file was copied to \filename{F:/public} on the
    intended controller;
    \item The file was copied completely and has not been renamed
    or modified;
    \item The controller was rebooted after the file was copied;
    \item The intended user and the resulting access permissions
    were verified after the restart.
\end{itemize}

If authorization is still not granted, consult the applicable
controller documentation or the responsible technical support
team before attempting the procedure again.

Once authorization has been confirmed, remove the file from
\filename{F:/public} immediately, regardless of whether any
additional troubleshooting is required.

\subsection{The reset file is not valid}

Check that:

\begin{itemize}
    \item The entered MAC address corresponds to the card;
    \item The MAC address format is correct;
    \item The card's date and time are correct;
    \item The file has not passed its expiration time.
\end{itemize}

\subsection{Firmware upload fails}

Check that:

\begin{itemize}
    \item The selected firmware is compatible with the card;
    \item The file was processed using
    \button{Firmware Checksum};
    \item The \filename{CHECKSUM.ENC} file was used;
    \item The operation was started from the
    \button{Upload Firmware} option.
\end{itemize}

\end{document}
